%---------------------------------------------------------%
%   CHAPTER 1: INTRODUCTION                                %
%---------------------------------------------------------%
\chapter{Introduction}

\section{Purpose}

This guideline is a practical reference for preparing the Projek Sarjana Muda (PSM) report at Universiti Malaysia Pahang (UMP). It combines two sources:

\begin{enumerate}
    \item the \textbf{UMP Thesis Guide} (November 2019), which defines the technical and formatting specifications; and
    \item the \textbf{PSM content guideline}, which defines what should be written in each chapter of the report.
\end{enumerate}

The formatting rules in this document follow the \emph{Publication Manual of the American Psychological Association}, Sixth Edition (APA 6th), as required by the UMP Thesis Guide \citep{umpthesis2019}. Where a rule can be produced with \LaTeX{}, this document shows the exact command or setting, based on the \texttt{umpthesis.cls} class \citep{umpthesiscls}.

\section{Definition of Terms}

\begin{description}
    \item[Typeface or font] All styles such as bold and italics, and the acceptable fonts such as Times New Roman, Arial, and Courier.
    \item[Text] The main body of the report. It excludes the cover page, preliminaries, tables, figures, references, and appendices.
    \item[Preliminaries] The front section of the report: title page, declaration pages, dedication (optional), acknowledgements, abstrak, abstract, table of contents, list of tables, list of figures, and list of abbreviations or symbols.
    \item[Table] Verbal, numeric, or graphical information arranged systematically in rows and columns.
    \item[Figure] Any information that is not tabulated: illustrations, graphs, charts, diagrams, photographs, or maps.
    \item[Equation] A mathematical or chemical expression that is part of the text. Equations must be centrally aligned.
    \item[Language] The report may be written in English or Bahasa Melayu. Other languages require approval from the University Senate.
    \item[Appendix] Supplementary material that is relevant to the text but would disrupt the flow if included in the main body.
\end{description}

\section{How to Use This Guideline}

\begin{itemize}
    \item \textbf{Chapter 2} lists the formatting rules and the matching \LaTeX{} setting.
    \item \textbf{Chapter 3} lists the required content of each PSM report chapter.
    \item \textbf{Chapter 4} gives step-by-step instructions for compiling a report with \texttt{umpthesis.cls}.
    \item \textbf{Chapter 5} provides a submission checklist and documents known differences between the class and the 2019 guide.
\end{itemize}
